\documentclass[11pt,a4paper]{article}
\usepackage[margin=1in]{geometry}
\usepackage{amsmath,amssymb,amsthm}
\usepackage{mathtools}
\usepackage{hyperref}
\usepackage{enumitem}
\usepackage{booktabs}
\usepackage{array}

\theoremstyle{plain}
\newtheorem{theorem}{Theorem}[section]
\newtheorem{proposition}[theorem]{Proposition}
\newtheorem{lemma}[theorem]{Lemma}
\newtheorem{corollary}[theorem]{Corollary}
\newtheorem{conjecture}[theorem]{Conjecture}

\theoremstyle{definition}
\newtheorem{definition}[theorem]{Definition}
\newtheorem{condition}[theorem]{Condition}

\theoremstyle{remark}
\newtheorem{remark}[theorem]{Remark}
\newtheorem{example}[theorem]{Example}

\newcommand{\R}{\mathbb{R}}
\newcommand{\Z}{\mathbb{Z}}
\newcommand{\T}{\mathbb{T}}
\newcommand{\N}{\mathbb{N}}
\newcommand{\E}{\mathbb{E}}
\newcommand{\Var}{\operatorname{Var}}
\newcommand{\Ent}{\operatorname{Ent}}
\newcommand{\tr}{\operatorname{tr}}
\newcommand{\Tr}{\operatorname{Tr}}
\newcommand{\Hess}{\operatorname{Hess}}
\newcommand{\Ric}{\operatorname{Ric}}
\newcommand{\gap}{\operatorname{gap}}
\newcommand{\ad}{\operatorname{ad}}
\newcommand{\Ad}{\operatorname{Ad}}
\newcommand{\SU}{\mathrm{SU}}
\newcommand{\SO}{\mathrm{SO}}
\newcommand{\U}{\mathrm{U}}
\newcommand{\fg}{\mathfrak{g}}
\newcommand{\su}{\mathfrak{su}}
\newcommand{\cA}{\mathcal{A}}
\newcommand{\cB}{\mathcal{B}}
\newcommand{\cC}{\mathcal{C}}
\newcommand{\cE}{\mathcal{E}}
\newcommand{\cF}{\mathcal{F}}
\newcommand{\cG}{\mathcal{G}}
\newcommand{\cH}{\mathcal{H}}
\newcommand{\cK}{\mathcal{K}}
\newcommand{\cL}{\mathcal{L}}
\newcommand{\cO}{\mathcal{O}}
\newcommand{\cP}{\mathcal{P}}
\newcommand{\cR}{\mathcal{R}}
\newcommand{\cS}{\mathcal{S}}
\newcommand{\cT}{\mathcal{T}}
\newcommand{\cV}{\mathcal{V}}
\newcommand{\cW}{\mathcal{W}}
\newcommand{\cZ}{\mathcal{Z}}
\newcommand{\bA}{\mathbf{A}}
\newcommand{\bB}{\mathbf{B}}
\newcommand{\bU}{\mathbf{U}}

\title{UV Stability for Non-Abelian Gauge Theories\\
  via the Dimock--Balaban Renormalization Group:\\
  A Roadmap from $\phi^4_3$ and QED$_3$ to YM$_{3,4}$}
\author{R.\ da Costa Fontes\\ \textit{Quijotic Research}\\
\texttt{ramirofontes@gmail.com}}
\date{February 2026}

\begin{document}
\maketitle

\begin{abstract}
Dimock~\cite{Dimock13a,Dimock13b,Dimock14} gave a complete,
self-contained treatment of the ultraviolet problem for
$\phi^4_3$ using Balaban's renormalization group method,
and subsequently extended this to QED in $d=3$
in a series of papers~\cite{Dimock02,Dimock04,Dimock15a,Dimock15b,Dimock18,Dimock22}.
We present a detailed roadmap for extending Dimock's
framework to \emph{non-Abelian} gauge theories:
first to pure $\SU(N)$ Yang--Mills in $d=3$
(super-renormalizable), then to $d=4$ (just-renormalizable,
asymptotically free).

For each step of the Balaban RG procedure ---
block averaging, gauge fixing, variational problem,
fluctuation integration, small/large field decomposition,
and polymer expansion --- we identify precisely which
estimates from Dimock's work carry over to the
non-Abelian setting and which require genuinely new analysis.
We formulate the UV stability bound
as a \emph{Hessian condition} compatible with the
Bauerschmidt--Bodineau--Dagallier (BBD) multiscale
Bakry--\'Emery criterion~\cite{BBD24}, establishing
a direct interface between the Balaban UV program
and the author's mass gap framework~\cite{FontesBBD}.

The paper serves three purposes: (i)~a readable introduction
to the Dimock--Balaban method for researchers approaching
from the constructive QFT or probability communities;
(ii)~a precise statement of the estimates needed to
complete the non-Abelian UV program;
(iii)~a bridge connecting the UV renormalization group
to the BBD log-Sobolev framework, showing that the
Balaban effective action bounds are precisely the
missing UV input for the continuum Yang--Mills mass gap.
\end{abstract}

\noindent\textbf{MSC 2020:} 81T13, 81T25, 82B28, 60J25.

%% ====================================================================
\section{Introduction}\label{sec:intro}
%% ====================================================================

\subsection{The UV problem for Yang--Mills}

The construction of a quantum Yang--Mills theory satisfying
the Wightman or Osterwalder--Schrader axioms requires control
of two distinct problems: the \emph{ultraviolet} problem
(existence of the continuum limit as lattice spacing $a \to 0$)
and the \emph{infrared} problem (existence of the infinite-volume
limit and demonstration of a mass gap). In our companion
papers~\cite{FontesGap,FontesUniform,FontesBBD,FontesEuclidean},
we have developed a framework based on the Bakry--\'Emery
curvature condition that reduces the mass gap problem to
three scale-separated inputs:
\begin{equation}\label{eq:three_regimes}
  \text{IR curvature} \quad + \quad
  \text{crossover (block gap)} \quad + \quad
  \text{UV stability.}
\end{equation}
The first two inputs have been established for
3D Yang--Mills~\cite{FontesGap,FontesUniform}.
The UV stability bound is the focus of the present paper.

\subsection{Balaban's program and Dimock's modernization}

Balaban's monumental series of papers~\cite{Bal84a,Bal84b,Bal85a,Bal85b,Bal85c,Bal85d,Bal85e,Bal87,Bal88a,Bal88b,Bal89}
established UV stability for 3D lattice pure gauge theories
and developed the renormalization group machinery for 4D
gauge theories. However, these papers --- totaling approximately
800 pages --- are extremely dense and notation-heavy, with
key arguments distributed across many publications.

Dimock undertook a systematic modernization of Balaban's approach:
\begin{enumerate}[label=(\roman*)]
  \item \emph{$\phi^4_3$} (scalar, super-renormalizable):
    Three papers~\cite{Dimock13a,Dimock13b,Dimock14}
    providing a complete, self-contained proof of UV stability.
    Total: $\sim$200 pages, ``readable line by
    line''~\cite{Adhikari}.
  \item \emph{QED$_3$} (Abelian gauge + fermions,
    super-renormalizable): Seven papers~\cite{Dimock02,Dimock04,Dimock15a,Dimock15b,Dimock18,Dimock20,Dimock22}
    extending to $\U(1)$ gauge theory.
    Final result: UV stability bound for QED$_3$
    at weak coupling in a fixed finite volume.
\end{enumerate}
The natural next step is pure $\SU(N)$ Yang--Mills in $d=3$,
then $d=4$. This has \emph{not} been done; the purpose
of this paper is to lay out precisely what is needed.

\subsection{Structure of the paper}

Section~\ref{sec:balaban_rg} reviews the Balaban RG procedure
in Dimock's clean notation.
Section~\ref{sec:abelian_to_nonabelian} identifies the
obstacles in passing from $\U(1)$ to $\SU(N)$.
Section~\ref{sec:ym3} formulates the 3D non-Abelian program
and states the required estimates.
Section~\ref{sec:ym4} treats the 4D case with its
additional running coupling constant.
Section~\ref{sec:bbd_interface} connects the effective
action bounds to the BBD Hessian condition.
Section~\ref{sec:summary} summarizes the status
and prioritizes the open problems.

%% ====================================================================
\section{The Balaban RG Procedure}\label{sec:balaban_rg}
%% ====================================================================

We follow Dimock's notation~\cite{Dimock13a,Dimock22}.
The theory lives on a toroidal lattice
$\T_\eta^d = (\eta \Z / L^M \Z)^d$ with lattice spacing
$\eta = L^{-N}$ and physical volume $L^M$, where $L$ is a
fixed large odd integer and $M, N$ are large positive integers.
The continuum limit is $N \to \infty$; the
infinite-volume limit is $M \to \infty$.

\subsection{Lattice gauge fields}

A \emph{lattice gauge field} $U$ assigns to each oriented
bond $b = \langle x, x + \eta e_\mu \rangle$ an element
$U(b) \in G$, with $U(-b) = U(b)^{-1}$.
For a plaquette $p$ with boundary bonds $b_1, \ldots, b_4$,
the plaquette holonomy is
\begin{equation}
  U(p) = U(b_1) U(b_2) U(b_3) U(b_4).
\end{equation}
The \emph{Wilson action} is
\begin{equation}\label{eq:wilson_action}
  S(U) = \frac{1}{g^2} \sum_p \eta^{d-4}
  \left(1 - \frac{1}{\dim R}\,\text{Re}\,\tr_R\, U(p)\right),
\end{equation}
where $R$ is the fundamental representation and the sum is
over unoriented plaquettes. The factor $\eta^{d-4}$ ensures
the correct continuum scaling.

A \emph{lattice gauge transformation} $u: x \mapsto u(x) \in G$
acts by $U(b) \mapsto u(x) U(b) u(y)^{-1}$ for $b = \langle x,y\rangle$.
The partition function is
\begin{equation}
  \cZ = \int \prod_b dU(b)\; e^{-S(U)},
\end{equation}
where $dU(b)$ is the Haar measure on $G$.

\subsection{Block averaging}\label{sec:block_avg}

The key operation in Balaban's RG is
\emph{gauge-covariant block averaging}.
Starting on the fine lattice $\T_\eta^d$, one averages
to a coarser lattice $\T_{L\eta}^d$ with spacing $L\eta$.

For a bond $B = \langle X, X + L\eta\, e_\mu \rangle$
on the coarse lattice, the \emph{block average} $Q(U, B)$
is defined via a gauge-covariant sum over paths
from $X$ to $X + L\eta\, e_\mu$ on the fine lattice.
Specifically, let $\Gamma(B)$ denote the set of
shortest lattice paths connecting the endpoints of $B$
through the fine lattice. Then
\begin{equation}\label{eq:block_avg}
  Q(U, B) = \frac{1}{|\Gamma(B)|} \sum_{\gamma \in \Gamma(B)}
  U(\gamma),
\end{equation}
where $U(\gamma) = U(b_1) \cdots U(b_n)$ is the ordered
product along the path. The crucial property is:
\begin{equation}
  Q(U^u, B) = u(X)\, Q(U, B)\, u(X + L\eta\, e_\mu)^{-1},
\end{equation}
i.e., $Q$ transforms covariantly under gauge transformations.

\begin{remark}
  For $\U(1)$, $Q(U,B)$ takes values in $\mathbb{C}$
  and lies near the unit circle for small fields.
  For $\SU(N)$, $Q(U,B)$ is an $N \times N$ matrix
  near $\SU(N)$ but \emph{not necessarily in $\SU(N)$}.
  This is the first non-Abelian subtlety: one must
  project back to the group, typically by normalizing
  $Q \to Q / |\det Q|^{1/N}$ or using the polar
  decomposition. See Section~\ref{sec:nonabelian_avg}.
\end{remark}

\subsection{The RG step}\label{sec:rg_step}

One RG step transforms the density
$\rho_k(\Phi_k)$ on the $k$-th scale lattice into a density
$\rho_{k+1}(\Phi_{k+1})$ on the $(k+1)$-th scale lattice.
The procedure has five stages:

\medskip\noindent
\textbf{Stage 1: Constrained integration.}
Fix the block averages $\Phi_{k+1}$ and integrate
over configurations $\Phi_k$ on the finer lattice
with the constraint $Q(\Phi_k) = \Phi_{k+1}$:
\begin{equation}\label{eq:constrained}
  \rho_{k+1}(\Phi_{k+1}) = \int
  \delta(Q(\Phi_k) - \Phi_{k+1})\;
  \rho_k(\Phi_k)\; d\Phi_k.
\end{equation}

\noindent
\textbf{Stage 2: Background field (variational problem).}
Find the \emph{minimizer} $\Phi_k^{\min}$ of the action
subject to $Q(\Phi_k^{\min}) = \Phi_{k+1}$.
This background field satisfies the lattice
Euler--Lagrange equations with Lagrange multipliers:
\begin{equation}\label{eq:variational}
  \frac{\partial S}{\partial \Phi_k}
  \bigg|_{\Phi_k^{\min}} =
  Q^* \lambda,
\end{equation}
where $Q^*$ is the adjoint of the linearized averaging
operator. The minimizer is unique up to gauge transformations
when $\Phi_{k+1}$ is in a ``small field'' region.

\medskip\noindent
\textbf{Stage 3: Fluctuation field.}
Decompose $\Phi_k = \Phi_k^{\min} \cdot \phi$
where $\phi$ is the fluctuation. For gauge fields
this is multiplicative: $U(b) = U^{\min}(b)\, e^{i\eta A(b)}$
where $A(b) \in \fg$ is the Lie-algebra-valued fluctuation.

\medskip\noindent
\textbf{Stage 4: Small/large field decomposition.}
Split the lattice into regions where the fluctuation
field is ``small'' (controlled by perturbation theory)
and ``large'' (controlled by positivity of the action):
\begin{equation}
  \rho_{k+1} = \rho_{k+1}^{\text{small}} +
  \rho_{k+1}^{\text{large}}.
\end{equation}
In the small field region, expand the action to
second order in $\phi$ and perform Gaussian integration.
This produces the \emph{effective action} at scale $k+1$.

\medskip\noindent
\textbf{Stage 5: Gauge fixing.}
To define the Gaussian integral over $\phi$, one must
fix the gauge. Balaban uses \emph{covariant axial gauge}
within blocks, which is adapted to the block structure
and maintains locality. Dimock~\cite{Dimock15a} analyzes
this in detail for the Abelian case.

\subsection{After $N$ steps}

After $N$ RG steps, the fine lattice $\T_{L^{-N}}^d$
has been coarsened to the unit lattice $\T_1^d$.
The partition function takes the form
\begin{equation}\label{eq:after_N}
  \cZ = \int d\Phi_N \;
  e^{-S_N^{\text{eff}}(\Phi_N)}\;
  K_N(\Phi_N),
\end{equation}
where $S_N^{\text{eff}}$ is the effective action
incorporating all counterterms, and $K_N$ contains
the ``remainder'' terms from the cluster/polymer expansion.
The \textbf{UV stability bound} states:
\begin{equation}\label{eq:uv_stability}
  c_1 \le \cZ \le c_2,
\end{equation}
with constants $c_1, c_2 > 0$ \emph{uniform in $N$}.
This is the main result of Balaban~\cite{Bal85e}
for 3D and the goal for 4D.

%% ====================================================================
\section{From Abelian to Non-Abelian:
  The Six Obstacles}\label{sec:abelian_to_nonabelian}
%% ====================================================================

We now systematically compare Dimock's QED$_3$
treatment with what is needed for $\SU(N)$ YM.
We identify six specific obstacles, ranging from
minor notation changes to genuinely hard analysis.

\subsection{Obstacle 1: Block averaging (moderate)}
\label{sec:nonabelian_avg}

For $\U(1)$, the block average~\eqref{eq:block_avg} produces
a complex number near the unit circle. For $\SU(N)$,
the average of group elements is a matrix
that is \emph{not} in $\SU(N)$ but lies in
$\text{GL}(N, \mathbb{C})$ near $\SU(N)$.

Balaban's solution~\cite{Bal85a}: define the averaging
operation in terms of Lie algebra components.
For small fields, write $U(b) = e^{i\eta A_\mu(x)}$
and define the block average via
\begin{equation}\label{eq:lie_avg}
  (QA)_\mu(X) = \frac{1}{L}
  \sum_{j=0}^{L-1} \left(
  \text{parallel transport to } X\right)
  A_\mu(X + j\eta e_\mu),
\end{equation}
where the parallel transport uses the background field.
This is gauge-covariant but requires the background
field to be specified first, creating a circular
dependence that must be resolved iteratively.

\begin{remark}[Dimock's approach]
  In~\cite{Dimock22}, Dimock defines the averaging operator
  for $\U(1)$ gauge fields as a simple sum:
  $(QA)(B) = \sum_{b \subset B} A(b)$.
  For $\SU(N)$, this formula works only for
  infinitesimal fields. The full non-Abelian definition
  requires the parallel transport and is the source of
  significant additional complexity in the estimates.
\end{remark}

\textbf{Difficulty:} Moderate. Balaban~\cite{Bal85a}
defined and analyzed the non-Abelian averaging operation.
The key estimates (regularity, analyticity on small field
configurations) are proved in~\cite{Bal85a} and need
only to be translated into Dimock's cleaner notation.

\subsection{Obstacle 2: Gauge fixing (moderate)}

In QED$_3$, Dimock uses the \emph{Landau gauge}
$\partial_\mu A_\mu = 0$, which is linear and
has well-understood properties. For non-Abelian
theories, one must use a gauge condition that is
compatible with the block structure.

Balaban uses \emph{covariant axial gauge within blocks},
which is adapted to the hierarchical structure of the RG.
Dimock~\cite{Dimock15a} already studied covariant axial
gauge in the Abelian setting. The non-Abelian version
requires solving the gauge-fixing condition
\begin{equation}\label{eq:axial_gauge}
  A_d(x_1, \ldots, x_{d-1}, x_d^0) = 0
\end{equation}
along ``axial trees'' within each block, where $x_d^0$
is a fixed reference value. This is well-defined for
small fields but requires the implicit function theorem
on the gauge group.

\textbf{Difficulty:} Moderate. The main issue is
ensuring that the Faddeev--Popov determinant is
non-degenerate on the small field region.
For $\U(1)$ this is automatic; for $\SU(N)$ it
requires a Gribov-type argument restricted to the
small field domain.

\subsection{Obstacle 3: Variational problem (hard)}

The \emph{variational problem} --- finding the
minimizer $U^{\min}$ of the Wilson action
subject to fixed block averages --- is the technical
heart of Balaban's approach.

For scalars, the minimizer satisfies a linear
(Laplace) equation. For $\U(1)$ gauge theory,
the linearized problem is still a Laplace equation
on 1-forms. For $\SU(N)$, the Euler--Lagrange
equation is a \emph{nonlinear} PDE on the lattice:
\begin{equation}\label{eq:euler_lagrange}
  d_U^* F_U + Q^* \lambda = 0,
\end{equation}
where $d_U^* F_U = \sum_\mu (D_\mu^{U})^* F_{\mu\nu}^U$
is the lattice Yang--Mills equation with covariant derivative
$D_\mu^U = \nabla_\mu + [A_\mu, \cdot]$, and $\lambda$
is the Lagrange multiplier for the averaging constraint.

\textbf{What Balaban proved:}
\begin{enumerate}[label=(\alph*)]
  \item Existence and uniqueness (modulo gauge) of the
    minimizer on the small field region~\cite{Bal85d}.
  \item Regularity: the minimizer is analytic in the
    block-averaged field $\Phi_{k+1}$~\cite{Bal85d}.
  \item Decay: the dependence on boundary data
    decays exponentially~\cite{Bal85d}.
\end{enumerate}
All three properties require the Banach contraction
mapping theorem applied to the nonlinear equation,
using the smallness of $g^2 \eta^{(d-4)/2}$ (or
$g^2(k)$ in $d=4$).

\textbf{Difficulty:} Hard. This is the most labor-intensive
part of the non-Abelian extension. The nonlinearity of
the Euler--Lagrange equation requires careful control
of commutator terms. The recent paper of
Bałaban and Imbrie~\cite{BI_sigma} studies the
analogous variational problem for non-linear sigma models
and provides a template.

\subsection{Obstacle 4: Fluctuation propagator (hard)}

After expanding around the minimizer, the Gaussian
integral over fluctuation fields $\phi$ involves a
propagator $C_k$ that depends on the \emph{background
field} $U^{\min}$:
\begin{equation}\label{eq:propagator}
  C_k = \left(-D_U^2 + Q^* a Q + \text{gauge fixing}
  \right)^{-1},
\end{equation}
where $D_U^2$ is the gauge-covariant Laplacian in
the background $U^{\min}$, and $Q^* a Q$ is the
block averaging constraint.

For $\U(1)$, $D_U^2 = \Delta + i[\nabla A, \cdot]$
reduces to a scalar Laplacian with lower-order terms.
For $\SU(N)$, the covariant Laplacian involves
commutator terms $[A_\mu, [A_\mu, \cdot]]$ that
create Lie-algebra-valued matrix structure.

\textbf{What is needed:}
\begin{enumerate}[label=(\alph*)]
  \item Existence of $C_k$ (invertibility of the operator).
  \item Exponential decay: $|C_k(x,y)| \le C e^{-m|x-y|}$.
  \item Random walk expansion: representation of $C_k$
    as a convergent sum over lattice paths.
  \item Analyticity in the background field.
\end{enumerate}
Balaban proved all four for non-Abelian backgrounds
in~\cite{Bal85b}. The random walk expansion is the
most delicate, relying on a resolvent identity and
Neumann series in the commutator terms.

\textbf{Difficulty:} Hard but with a clear path.
Dimock~\cite{Dimock22} has all four properties for
$\U(1)$ backgrounds. The extension to $\SU(N)$ requires
controlling the additional commutator terms, which are
proportional to $g^2$ and hence small in the UV regime.

\subsection{Obstacle 5: Polymer/cluster expansion
  (moderate--hard)}

The effective action after Gaussian integration is
represented as a \emph{polymer expansion}:
\begin{equation}\label{eq:polymer}
  e^{-S_{k+1}^{\text{eff}}(\Phi_{k+1})} =
  \sum_{\{X_i\}} \prod_i K(X_i, \Phi_{k+1}),
\end{equation}
where the sum is over collections of disjoint ``polymers''
(connected subsets of the coarse lattice) and
$K(X, \Phi_{k+1})$ is the activity of polymer $X$.

The convergence of this expansion requires:
\begin{enumerate}[label=(\alph*)]
  \item Exponential decay of $K(X, \Phi_{k+1})$
    in the size of $X$.
  \item Analyticity of $K$ in $\Phi_{k+1}$ on the
    small field region.
  \item Stable bounds under iteration of RG steps.
\end{enumerate}
For scalars, the polymer expansion is standard
(Brydges~\cite{Brydges86}). For gauge theories,
gauge invariance imposes additional structure:
each polymer activity must be gauge-invariant,
and the expansion must respect this.

\textbf{Difficulty:} Moderate--hard. The polymer
expansion itself is well-understood technology.
The additional difficulty is tracking the Lie algebra
structure through the expansion, ensuring that
commutator terms do not destroy the convergence
estimates.

\subsection{Obstacle 6: Large field bounds (moderate)}

The large field region $\Omega_k^c$ (where the
fluctuation field exceeds a threshold) is controlled
not by perturbation theory but by the \emph{positivity}
of the Wilson action:
\begin{equation}\label{eq:large_field}
  e^{-S(U)} \le e^{-c|\Omega_k^c| / \eta^d}
\end{equation}
for some constant $c > 0$ depending on $g^2$.
For non-Abelian gauge theories, the Wilson action
satisfies $S(U) \ge 0$ with equality only for
pure gauge configurations, and the positivity bound
is gauge-group-independent.

\textbf{Difficulty:} Moderate. Dimock's large field
analysis~\cite{Dimock13b} for $\phi^4_3$ carries over
with minor modifications. The key bound uses only
$1 - \text{Re}\,\tr U(p)/\dim R \ge 0$, which holds
for any compact group.

%% ====================================================================
\section{The 3D Non-Abelian Program}\label{sec:ym3}
%% ====================================================================

In $d = 3$, the Yang--Mills coupling $g^2$ has
mass dimension $[g^2] = 1$. The theory is
\emph{super-renormalizable}: the only divergent
counterterm is the vacuum energy density $\epsilon_k$
(cosmological constant). There is no coupling constant
renormalization and no mass renormalization at the
perturbative level. Non-perturbatively, a mass counterterm
of order $g^4$ appears~\cite{CS25}, but this is finite.

\subsection{Setup}

Fix $G = \SU(N)$ and $g^2 > 0$. The lattice
theory lives on $\T_\eta^3$ with $\eta = L^{-N}$
and inverse coupling $\beta = 1/(g^2 \eta)$.
As $\eta \to 0$ (equivalently $N \to \infty$),
$\beta \to \infty$ and the theory becomes weakly
coupled on the lattice scale.

The partition function is
\begin{equation}
  \cZ_{M,N} = \int \prod_b dU(b) \;
  e^{-\beta \sum_p (1 - \text{Re}\,\tr_R\,U(p)/N)}.
\end{equation}

\subsection{The effective action after $k$ steps}

After $k$ RG steps ($0 \le k \le N$), the effective
density on the scale-$k$ lattice $\T_{L^{-N+k}}^3$
takes the form
\begin{equation}\label{eq:effective_density_3d}
  \rho_k(\Phi_k) = e^{-\beta_k S_k(\Phi_k)
  - \epsilon_k |\Lambda_k| - V_k(\Phi_k)}
  \cdot K_k(\Phi_k),
\end{equation}
where:
\begin{itemize}[nosep]
  \item $\beta_k = L^k \beta$ is the rescaled inverse coupling
    (grows with $k$ due to super-renormalizability),
  \item $S_k$ is the Wilson action on $\T_{L^{-N+k}}^3$,
  \item $\epsilon_k$ is the vacuum energy counterterm,
  \item $V_k(\Phi_k) = O(\beta_k^{-1/2})$ is the
    ``relevant perturbation'' (small correction to the action),
  \item $K_k$ is the polymer expansion remainder
    with $\|K_k - 1\| = O(\beta_k^{-1/2})$.
\end{itemize}

\begin{remark}[Comparison with QED$_3$]
  In Dimock's QED$_3$ treatment~\cite{Dimock22},
  the effective density has the same structure
  with additional fermion determinant factors.
  For pure YM$_3$, the absence of fermions
  \emph{simplifies} the analysis: there is no
  fermion propagator, no fermion mass renormalization,
  and no Grassmann integration. The non-Abelian
  gauge field self-interaction is the only new ingredient.
\end{remark}

\subsection{Required estimates}

We now state the estimates needed for 3D non-Abelian
UV stability, organized by the six obstacles of
Section~\ref{sec:abelian_to_nonabelian}.

\begin{condition}[Non-Abelian block averaging]
\label{cond:avg}
For the $\SU(N)$ block averaging operator
$Q$ defined via~\eqref{eq:lie_avg} with background
field $U^{\min}$:
\begin{enumerate}[label=(\alph*),nosep]
  \item $Q$ is analytic on the small field region
    $\{\|F_U\|_\infty \le p_k\}$ with
    $p_k = O(\beta_k^{-1/4})$;
  \item the linearization $dQ$ satisfies
    $\|dQ - dQ_0\| \le C g^2 \eta \|A\|$ where
    $dQ_0$ is the free (Abelian) linearization;
  \item $Q$ transforms covariantly:
    $Q(U^u) = u_X \cdot Q(U) \cdot u_Y^{-1}$
    for any gauge transformation $u$.
\end{enumerate}
\end{condition}

\begin{condition}[Non-Abelian variational problem]
\label{cond:var}
For each block-averaged field $\Phi_{k+1}$ in the
small field region:
\begin{enumerate}[label=(\alph*),nosep]
  \item There exists a unique (modulo gauge) minimizer
    $U^{\min}$ of $S_k$ subject to $Q(U) = \Phi_{k+1}$;
  \item $U^{\min}$ is analytic in $\Phi_{k+1}$
    on the small field domain;
  \item the minimizer satisfies regularity:
    $\|F_{U^{\min}}\|_\infty \le C \|F_{\Phi_{k+1}}\|_\infty$;
  \item the dependence on data at distance $r$
    decays as $e^{-mr}$ with $m = O(\sqrt{\beta_k})$.
\end{enumerate}
\end{condition}

\begin{condition}[Non-Abelian fluctuation propagator]
\label{cond:prop}
The gauge-covariant propagator~\eqref{eq:propagator}
with $\SU(N)$ background satisfies:
\begin{enumerate}[label=(\alph*),nosep]
  \item $C_k$ exists and is bounded:
    $\|C_k\| \le C_0$ uniformly in the small field
    background;
  \item exponential decay: $|C_k(x,y)| \le C_0
    e^{-m_0 |x-y|/\eta_k}$ with $m_0 > 0$
    independent of $k$;
  \item the random walk expansion converges
    with $\SU(N)$ commutator corrections
    bounded by $C g^2 \eta_k$;
  \item $C_k$ is analytic in the background field
    on the small field domain.
\end{enumerate}
\end{condition}

\begin{condition}[Polymer expansion convergence]
\label{cond:poly}
The polymer activities $K(X, \Phi_{k+1})$
satisfy:
\begin{enumerate}[label=(\alph*),nosep]
  \item $|K(X, \Phi_{k+1})| \le e^{-c_0 |X|}$
    for a constant $c_0$ independent of $k$;
  \item gauge invariance: $K(X, \Phi_{k+1}^u) =
    K(X, \Phi_{k+1})$;
  \item stability: the estimates are preserved
    under $N \to N-1$ RG steps.
\end{enumerate}
\end{condition}

\begin{theorem}[3D non-Abelian UV stability ---
  conditional]\label{thm:3d_uv}
Assume Conditions~\ref{cond:avg}--\ref{cond:poly}
hold for $G = \SU(N)$ in $d = 3$ with coupling
$g^2 > 0$. Then the partition function
satisfies the UV stability bound:
\begin{equation}
  e^{-C_1 L^{3M}} \le \cZ_{M,N}
  \le e^{C_2 L^{3M}}
\end{equation}
uniformly in $N$ (lattice spacing $\eta = L^{-N}$),
for constants $C_1, C_2$ depending on $g^2$ and $L$
but not on $N$.
\end{theorem}

\begin{remark}[What is genuinely new vs.\ translation]
  Conditions~\ref{cond:avg}(a,c) and~\ref{cond:var}(a--d)
  are proved by Balaban~\cite{Bal85a,Bal85d} for $\SU(N)$
  in $d = 3,4$. The work needed is \emph{translation}
  into Dimock's framework: replacing Balaban's notation
  with Dimock's lattice calculus and verifying that
  all constants are tracked explicitly.
  Condition~\ref{cond:prop} requires a new calculation:
  the random walk expansion for the $\SU(N)$ propagator,
  extending Dimock's $\U(1)$ version~\cite{Dimock22}.
  The commutator terms create matrix-valued perturbations
  of the Laplacian that must be controlled.
\end{remark}

\subsection{Comparison table: QED$_3$ vs.\ YM$_3$}

\begin{table}[h]
\centering
\renewcommand{\arraystretch}{1.3}
\begin{tabular}{@{}p{4cm}ccp{4.5cm}@{}}
\toprule
\textbf{Component} & \textbf{QED$_3$} & \textbf{YM$_3$}
  & \textbf{New difficulty} \\
\midrule
Block averaging & Linear & Non-linear &
  Parallel transport, projection to $G$ \\
Gauge fixing & Landau & Covariant axial &
  Faddeev--Popov on $\SU(N)$ \\
Variational problem & Linear PDE & Nonlinear PDE &
  Contraction mapping on Lie algebra \\
Propagator & Scalar + perturbation &
  Matrix-valued & Commutator terms
  $[A_\mu, \cdot]$ \\
Polymer expansion & Standard & Standard +
  Lie structure & Track gauge invariance \\
Large field & $e^{-S} \ge 0$ & Same &
  None (group-independent) \\
Fermion determinant & Present & Absent &
  Simpler (no Grassmann) \\
Counterterms & $\epsilon_k, m_k^{\text{fermion}}$
  & $\epsilon_k$ only & Simpler \\
\bottomrule
\end{tabular}
\caption{Comparison of Dimock's QED$_3$ program
  with the proposed YM$_3$ extension.}
\label{tab:comparison}
\end{table}

\subsection{Estimate of effort}

Dimock's QED$_3$ program spans seven papers and
approximately 400 pages. We estimate the pure YM$_3$
extension as follows:
\begin{itemize}[nosep]
  \item Removal of fermion sector: saves $\sim$150 pages.
  \item Non-Abelian averaging + variational problem:
    adds $\sim$80 pages (translating Balaban).
  \item Non-Abelian propagator: adds $\sim$60 pages
    (new random walk expansion).
  \item Polymer expansion with Lie structure:
    adds $\sim$30 pages.
\end{itemize}
Net estimate: $\sim$220 pages, comparable to the
$\phi^4_3$ treatment. This is a realistic project
for 1--2 years of focused work.

%% ====================================================================
\section{The 4D Non-Abelian Program}\label{sec:ym4}
%% ====================================================================

The $d = 4$ case is qualitatively different.
The coupling constant $g^2$ is dimensionless,
asymptotic freedom produces a \emph{running coupling}
$g^2(k)$ that decreases logarithmically under the
RG flow, and there is a genuine coupling constant
renormalization.

\subsection{Running coupling and $\beta$-function}

After $k$ RG steps, the effective coupling satisfies
\begin{equation}\label{eq:running_coupling}
  \frac{1}{g^2(k)} = \frac{1}{g^2(0)} + b_0 k \log L
  + O(g^2(0)),
\end{equation}
where $b_0 = \frac{11N}{48\pi^2}$ for $\SU(N)$
is the one-loop $\beta$-function coefficient.
Asymptotic freedom means $g^2(k) \to 0$ as
$k \to 0$ (UV direction), and $g^2(k)$ grows
toward the IR.

The RG flow must be controlled for $0 \le k \le N$,
where $g^2(k)$ ranges from $g^2(0) \ll 1$
(the bare coupling at UV cutoff) to $g^2(N) \sim O(1)$
(the physical coupling at the lattice scale).

\subsection{Additional counterterms}

In $d = 4$, beyond the vacuum energy $\epsilon_k$,
one needs a \emph{coupling constant counterterm}
$\delta g^2(k)$ at each scale. The renormalization
condition is:
\begin{equation}\label{eq:renorm_condition}
  g^2_{\text{eff}}(k+1) = g^2_{\text{eff}}(k)
  - b_0 g^4(k) \log L + \delta g^2(k),
\end{equation}
where $\delta g^2(k) = O(g^6(k))$ is chosen to
maintain the flow on the ``renormalized trajectory.''

This is the content of Balaban's 4D papers~%
\cite{Bal87,Bal88a,Bal88b,Bal89}: constructing the
$\beta$-function non-perturbatively and showing that
the effective action at each scale has the correct
structure. The effective density becomes:
\begin{equation}\label{eq:effective_density_4d}
  \rho_k(\Phi_k) = e^{-\frac{1}{g^2(k)} S_k(\Phi_k)
  - \epsilon_k |\Lambda_k| - V_k(\Phi_k)}
  \cdot K_k(\Phi_k),
\end{equation}
with the crucial difference that $1/g^2(k)$ is now
$k$-dependent and the correction $V_k = O(g^2(k))$
must be tracked through the flow.

\subsection{Required estimates for 4D}

All four conditions from Section~\ref{sec:ym3}
apply, with modifications:

\begin{condition}[4D coupling constant flow]
\label{cond:4d}
The bare coupling $g^2(0)$ can be chosen as a
function of $N$ (the number of RG steps) such that:
\begin{enumerate}[label=(\alph*),nosep]
  \item $g^2(0) \sim \frac{1}{b_0 N \log L}$ as $N \to \infty$;
  \item the effective coupling $g^2(k)$ satisfies
    $c_1/(N-k) \le g^2(k) \le c_2/(N-k)$ for
    $0 \le k \le N - k_0$ (UV regime);
  \item the correction $V_k$ satisfies
    $\|V_k\|_\infty \le C g^4(k)$ uniformly;
  \item the polymer remainder satisfies
    $\|K_k - 1\| \le C g^2(k)$ uniformly.
\end{enumerate}
\end{condition}

\begin{condition}[4D effective action bounds]
\label{cond:4d_action}
At each scale $k$, the effective action
$\cS_k(\Phi_k) = g^{-2}(k) S_k + V_k$
satisfies the \emph{Hessian bound}:
\begin{equation}\label{eq:hessian_4d}
  \|\Hess \cS_k(\Phi_k) - \Hess_0\|_{\text{op}}
  \le C g^2(k)
\end{equation}
on the small field region, where $\Hess_0$ is the
Hessian of the free (Gaussian) action at scale $k$.
\end{condition}

\begin{theorem}[4D non-Abelian UV stability ---
  conditional]\label{thm:4d_uv}
Assume Conditions~\ref{cond:avg}--\ref{cond:poly}
(adapted to $d=4$ with running coupling)
and Conditions~\ref{cond:4d}--\ref{cond:4d_action}.
Then the partition function satisfies
\begin{equation}
  e^{-C_1 L^{4M}} \le \cZ_{M,N}
  \le e^{C_2 L^{4M}}
\end{equation}
uniformly in $N$, for constants depending on $L$
and the renormalized coupling $g^2_R = g^2(N)$.
\end{theorem}

This is the statement that Balaban's program
aims to prove in $d=4$. The key papers
are~\cite{Bal87,Bal88a,Bal88b,Bal89}, which establish
Condition~\ref{cond:4d}(a,b) in the small field
approximation and construct the $\beta$-function
to all orders in perturbation theory.
The \emph{completion} of this program ---
verifying convergence of the full expansion ---
remains open.

%% ====================================================================
\section{Interface with the BBD Framework}
\label{sec:bbd_interface}
%% ====================================================================

We now connect the Balaban UV effective action bounds
to the Bauerschmidt--Bodineau--Dagallier (BBD)
multiscale Bakry--\'Emery criterion~\cite{BBD24}.
This is the key new contribution: showing that the
Balaban program provides \emph{exactly} the UV input
needed for the mass gap proof.

\subsection{The BBD criterion (review)}

The BBD criterion (Theorem~3.6 of~\cite{BBD24})
provides a log-Sobolev inequality for a measure
$\mu = e^{-H} / Z$ from multiscale Hessian control.
We use the Polchinski flow formulation from~\cite{FontesBBD}.

Consider the Polchinski RG flow for
$V_t = -\log \E[e^{-H}|\cF_t]$, where $(\cF_t)$
is a filtration given by the RG decimation.
The BBD condition requires:
\begin{equation}\label{eq:bbd}
  \int_0^\infty \left\|
  \left(\Hess V_t - \gamma(t) \cdot \text{Id}
  \right)_+\right\|_\infty \; dt \le C_0 < \infty,
\end{equation}
where $\gamma(t)$ is the curvature of the reference
Gaussian measure at scale $t$ and $(\cdot)_+$
denotes the positive part.

\subsection{The UV contribution ($0 \le t \le t_1$)}

In the UV regime, the Polchinski flow corresponds
to the Balaban RG steps $k = 0, 1, \ldots, N - k_0$.
The connection is:
\begin{equation}
  t \leftrightarrow k, \qquad
  \eta_k = L^{k-N} \leftrightarrow e^{-t},
\end{equation}
and the potential $V_t$ at scale $t$ corresponds
to the Balaban effective correction $V_k$.

The Hessian bound~\eqref{eq:hessian_4d}
translates directly into the BBD integrand:
\begin{equation}\label{eq:uv_bbd}
  \left\|(\Hess V_t - \gamma(t))_+\right\|_\infty
  \le C g^4(t) \qquad \text{(3D)},
\end{equation}
\begin{equation}\label{eq:uv_bbd_4d}
  \left\|(\Hess V_t - \gamma(t))_+\right\|_\infty
  \le C g^2(t) \qquad \text{(4D)}.
\end{equation}

\begin{proposition}[UV contribution to BBD integral]
\label{prop:uv_bbd}
Assume the Balaban effective action bounds
(Condition~\ref{cond:4d_action} in 4D,
or the corresponding 3D bound $\|V_k\| \le C/\beta_k$).
Then the UV contribution to the BBD integral converges:
\begin{align}
  \text{3D:} &\quad
  \sum_{k=0}^{N} C / \beta_k =
  \sum_{k=0}^{N} C g^2 L^{k-N} =
  O(g^2) < \infty, \\
  \text{4D:} &\quad
  \sum_{k=0}^{N} C g^2(k) =
  \sum_{k=0}^{N} \frac{C}{b_0(N-k)\log L} =
  O(\log N).
\end{align}
In 3D, the sum converges uniformly in $N$
(geometric series with ratio $L^{-1}$).
In 4D, the sum grows logarithmically; however,
the BBD condition~\eqref{eq:bbd} requires only
finiteness for each $N$, and the mass gap
$m \ge c / \log(1/(g^2_R))$ is consistent
with the logarithmic growth.
\end{proposition}

\begin{remark}[The 3D case is complete]
  In 3D, the UV contribution is $O(g^2)$,
  the crossover contribution is controlled by the
  lattice block gap (our~\cite{FontesGap}), and the
  IR contribution is controlled by single-link
  Bakry--\'Emery curvature~\cite{FontesGap}.
  Therefore, \emph{assuming the Balaban 3D UV stability
  bound} (which is~\cite{Bal85e}, translated to
  Dimock notation), the full BBD condition is
  satisfied and the 3D mass gap is proved
  in the continuum limit.
\end{remark}

\subsection{The crossover and IR contributions (review)}

For completeness, the full BBD integral decomposes as:
\begin{align}
  \int_0^\infty (\cdots)\;dt &=
  \underbrace{\int_0^{t_1}}_{\text{UV: Balaban bounds}}
  + \underbrace{\int_{t_1}^{t_2}}_{\text{crossover:
    block gap}~\cite{FontesGap}}
  + \underbrace{\int_{t_2}^{\infty}}_{\text{IR:
    Bakry--\'Emery}~\cite{FontesGap}}. \label{eq:decomp}
\end{align}
The crossover scale $t_1$ is where $g^2(t_1) \sim O(1)$
(in 4D) or where $\beta_k \sim O(1)$ (in 3D).
The IR scale $t_2$ is where the effective theory
is a single-plaquette measure.

%% ====================================================================
\section{Summary and Open Problems}\label{sec:summary}
%% ====================================================================

\subsection{Status table}

\begin{table}[h]
\centering
\renewcommand{\arraystretch}{1.3}
\begin{tabular}{@{}p{5cm}ccp{3.5cm}@{}}
\toprule
\textbf{Component} & \textbf{3D} & \textbf{4D}
  & \textbf{Source} \\
\midrule
Lattice mass gap ($\gamma_0 > 0$) & \checkmark &
  \checkmark & \cite{FontesGap} \\
Uniform-in-$a$ gap & \checkmark & Conditional
  & \cite{FontesUniform} \\
Euclidean covariance & \checkmark & Conditional
  & \cite{FontesEuclidean} \\
BBD framework + decomposition & \checkmark &
  \checkmark & \cite{FontesBBD} \\
UV stability (Abelian, Dimock) & \checkmark$^*$
  & --- & \cite{Dimock22} \\
UV stability (non-Abelian, Balaban) & \checkmark
  & Partial & \cite{Bal85e,Bal87} \\
UV stability (non-Abelian, Dimock) & \textbf{Open}
  & \textbf{Open} & This paper \\
Coupling constant renorm.\ & N/A & Partial
  & \cite{Bal87,Bal88a} \\
Full 4D completion & N/A & \textbf{Open} & --- \\
\bottomrule
\end{tabular}
\caption{Status of the Yang--Mills construction program.
  $\checkmark^*$ denotes the Abelian case only.}
\label{tab:status}
\end{table}

\subsection{Priority ranking of open problems}

\begin{enumerate}[label=\textbf{P\arabic*.}]
  \item \textbf{Non-Abelian fluctuation propagator
    in 3D.} Extend Dimock's random walk expansion
    from $\U(1)$ to $\SU(N)$ backgrounds. This is the
    single hardest estimate for YM$_3$ and is
    \emph{necessary and sufficient} (combined with
    the translation of Balaban's other estimates) for
    completing the 3D construction.
    \emph{Estimated effort: 40--60 pages, 6 months.}

  \item \textbf{Non-Abelian variational problem
    in Dimock notation.} Translate Balaban's
    existence/uniqueness/regularity results
    for the minimizer from~\cite{Bal85d} into
    Dimock's framework. Template:
    Bałaban--Imbrie~\cite{BI_sigma} for nonlinear
    sigma models.
    \emph{Estimated effort: 60--80 pages, 8 months.}

  \item \textbf{4D $\beta$-function in BBD language.}
    Reformulate Balaban's non-perturbative
    $\beta$-function construction~\cite{Bal87,Bal88a}
    as a Hessian flow, showing that the BBD
    condition~\eqref{eq:bbd} is maintained under the
    coupling constant flow. This is the bridge
    between Balaban's 4D program and the mass gap.
    \emph{Estimated effort: 30--40 pages, 4 months.}

  \item \textbf{Completion of 4D large field bounds.}
    Balaban's treatment of the 4D large field
    region is incomplete in the published
    literature. Combine with the polymer expansion
    convergence to obtain the full 4D UV stability.
    \emph{Estimated effort: 80--100 pages, 12 months.}
\end{enumerate}

\subsection{The 3D milestone}

We emphasize that the 3D case is \emph{close to
complete}. The ingredients are:
\begin{enumerate}[label=(\roman*),nosep]
  \item Lattice mass gap: proved~\cite{FontesGap}.
  \item Uniform gap in $a$: proved~\cite{FontesUniform}.
  \item Euclidean covariance: proved~\cite{FontesEuclidean}.
  \item UV stability: Balaban~\cite{Bal85e} (non-Abelian,
    3D, published 1985, but in Balaban's notation).
  \item Translation to Dimock: open (this paper's
    roadmap), estimated $\sim$220 pages.
\end{enumerate}
If one accepts Balaban's 1985 result as stated,
the 3D Yang--Mills construction is complete modulo
the translation exercise. The value of the Dimock
modernization is not mathematical novelty but
\emph{verifiability}: producing a self-contained,
readable proof that can be checked by the community.

\subsection{Connection to Chevyrev--Shen}

The recent work of Chevyrev and Shen~\cite{CS25}
on uniqueness of gauge-covariant mass renormalization
in 3D provides an important complement to the
Dimock--Balaban approach. Their result ensures that
the continuum limit is \emph{independent of the
lattice regularization scheme}, which is precisely
what our UV universality proposition
in~\cite{FontesEuclidean} requires. The two programs
are synergistic: Chevyrev--Shen provides the
\emph{uniqueness} of the continuum theory, while
Dimock--Balaban provides the \emph{existence}
via controlled UV renormalization.

%% ====================================================================
\section*{Acknowledgments}
%% ====================================================================

The author thanks Claude (Anthropic) for assistance
in organizing the technical comparison between the
Abelian and non-Abelian programs. This work is part
of the Yang--Mills mass gap program at Quijotic Research.

\begin{thebibliography}{99}

\bibitem{Adhikari} A.~Adhikari, R.~Bauerschmidt, and
  T.~Helmuth, ``Spectral gap and log-Sobolev inequality
  for lattice gauge theories,'' in preparation (cited
  in BBD~\cite{BBD24}).

\bibitem{Bal84a} T.~Ba{\l}aban, ``Propagators and
  renormalization transformations for lattice gauge
  field theories~I,''
  Commun.\ Math.\ Phys.\ \textbf{95} (1984) 17--40.

\bibitem{Bal84b} T.~Ba{\l}aban, ``Propagators and
  renormalization transformations for lattice gauge
  field theories~II,''
  Commun.\ Math.\ Phys.\ \textbf{96} (1984) 223--250.

\bibitem{Bal85a} T.~Ba{\l}aban, ``Averaging operations
  for lattice gauge theories,''
  Commun.\ Math.\ Phys.\ \textbf{98} (1985) 17--51.

\bibitem{Bal85b} T.~Ba{\l}aban, ``Propagators for
  lattice gauge theories in a background field,''
  Commun.\ Math.\ Phys.\ \textbf{99} (1985) 389--434.

\bibitem{Bal85c} T.~Ba{\l}aban, ``Spaces of regular
  gauge field configurations on a lattice and gauge
  fixing conditions,''
  Commun.\ Math.\ Phys.\ \textbf{99} (1985) 75--102.

\bibitem{Bal85d} T.~Ba{\l}aban, ``The variational problem
  and background fields in renormalization group method
  for lattice gauge theories,''
  Commun.\ Math.\ Phys.\ \textbf{102} (1985) 277--309.

\bibitem{Bal85e} T.~Ba{\l}aban, ``Ultraviolet stability
  of three-dimensional lattice pure gauge field theories,''
  Commun.\ Math.\ Phys.\ \textbf{102} (1985) 255--275.

\bibitem{Bal87} T.~Ba{\l}aban, ``Renormalization group
  approach to lattice gauge field theories.~I,''
  Commun.\ Math.\ Phys.\ \textbf{109} (1987) 249--301.

\bibitem{Bal88a} T.~Ba{\l}aban, ``Large field
  renormalization.~I,''
  Commun.\ Math.\ Phys.\ \textbf{122} (1989) 175--202.

\bibitem{Bal88b} T.~Ba{\l}aban, ``Large field
  renormalization.~II,''
  Commun.\ Math.\ Phys.\ \textbf{122} (1989) 355--392.

\bibitem{Bal89} T.~Ba{\l}aban, ``A low temperature
  expansion for classical $N$-vector models,''
  Commun.\ Math.\ Phys.\ \textbf{167} (1995) 103--154.

\bibitem{BBD24} R.~Bauerschmidt, T.~Bodineau, and
  B.~Dagallier, ``Log-Sobolev inequality for the
  continuum sine-Gordon model,''
  Commun.\ Pure Appl.\ Math.\ (2024).

\bibitem{BI_sigma} T.~Ba{\l}aban, ``The variational
  problem for non-linear sigma models,''
  see also discussion in Adhikari et al.

\bibitem{Brydges86} D.C.~Brydges, ``A short course on
  cluster expansions,'' in \emph{Critical Phenomena,
  Random Systems, Gauge Theories} (Les Houches 1984),
  North-Holland, 1986, pp.~129--183.

\bibitem{CS25} I.~Chevyrev and H.~Shen, ``Gauge-covariant
  mass renormalization in 3D Yang--Mills--Higgs theory,''
  arXiv:2503.03060 (2025).

\bibitem{Dimock02} J.~Dimock, ``Quantum electrodynamics
  on the 3-torus~I,'' arXiv:math-ph/0210020 (2002).

\bibitem{Dimock04} J.~Dimock, ``Quantum electrodynamics
  on the 3-torus~II,'' arXiv:math-ph/0407063 (2004).

\bibitem{Dimock13a} J.~Dimock, ``The renormalization group
  according to Balaban~I.~Small fields,''
  Rev.\ Math.\ Phys.\ \textbf{25} (2013) 1330010.

\bibitem{Dimock13b} J.~Dimock, ``The renormalization group
  according to Balaban~II.~Large fields,''
  J.~Math.\ Phys.\ \textbf{54} (2013) 092301.

\bibitem{Dimock14} J.~Dimock, ``The renormalization group
  according to Balaban~III.~Convergence,''
  Ann.\ Henri Poincar\'e \textbf{15} (2014) 2133--2175.

\bibitem{Dimock15a} J.~Dimock, ``Covariant axial gauge,''
  Lett.\ Math.\ Phys.\ \textbf{105} (2015) 959--987.

\bibitem{Dimock15b} J.~Dimock, ``Nonperturbative
  renormalization of scalar QED in $d=3$,''
  J.~Math.\ Phys.\ \textbf{56} (2015) 102304.

\bibitem{Dimock18} J.~Dimock, ``Ultraviolet regularity
  for QED in $d=3$,''
  J.~Math.\ Phys.\ \textbf{59} (2018) 012301.

\bibitem{Dimock20} J.~Dimock, ``Propagators and minimizers
  for QED in $d=3$'' (2020), see~\cite{Dimock22}.

\bibitem{Dimock22} J.~Dimock, ``Ultraviolet stability
  for QED in $d=3$,''
  Ann.\ Henri Poincar\'e \textbf{23} (2022) 2685--2727.

\bibitem{FontesBBD} R.~da Costa Fontes, ``Toward the
  continuum Yang--Mills mass gap: the Polchinski flow,
  lattice block gaps, and the multiscale Bakry--\'Emery
  criterion,'' Quijotic Research (2026).

\bibitem{FontesEuclidean} R.~da Costa Fontes, ``Euclidean
  covariance of the 3D Yang--Mills continuum limit,''
  Quijotic Research (2026).

\bibitem{FontesGap} R.~da Costa Fontes, ``A mass gap for
  3D lattice Yang--Mills via Bakry--\'Emery curvature,''
  Quijotic Research (2025--2026).

\bibitem{FontesUniform} R.~da Costa Fontes, ``Uniform
  spectral gap for 3D Yang--Mills in the continuum
  limit,'' Quijotic Research (2026).

\end{thebibliography}

\end{document}
